\subsection{EXTENSION OF MAPPINGS TO PRODUCTS}

DEFINITION 2. Let \((\mathbf{X}_t)_{t \in I}, (\mathbf{Y}_t)_{t \in I}\) be two families of sets, and let \((g_t)_{t \in I}\) be a family of functions with the same index set I such that \(g_t\) is a mapping of \(\mathbf{X}_t\) into \(\mathbf{Y}_t\) for each \(t \in I\) . For each \(f \in \prod_{t \in I} \mathbf{X}_t\) let \(u_f\) be the graph of the function \(t = g_t(f(t))\) ( \(t \in I\) ), which is an element of \(\prod_{t \in I} \mathbf{Y}_t\) . The mapping \(f \to u_f\) of \(\prod_{t \in I} \mathbf{X}_t\) into \(\prod_{t \in I} \mathbf{Y}_t\) is called the canonical extension (or simply the extension) to products of the family of mappings \((g_t)_{t \in I}\) ; it is also sometimes called the product of the family of mappings \((g_t)_{t \in I}\) .

When the index notation is used, the product of the family \((g_{i})_{i\in I}\) is the function \((x_{i})_{i\in I}\to (g_{i}(x_{i}))_{i\in I}\) ; it is sometimes denoted by \((g_{i})_{i\in I}\) .

If \(\mathbf{I} = \{\alpha, \beta\}\) , where \(\alpha\) and \(\beta\) are distinct, the extension to products of the family of mappings \((g_t)_{t \in \mathbf{I}}\) is just \(\psi \circ (g_\alpha \times g_\beta) \circ \varphi\) , where \(\varphi\) denotes the canonical mapping of \(\prod_{i \in \mathbf{I}} \mathbf{X}_i\) onto \(\mathbf{X}_\alpha \times \mathbf{X}_\beta\) (no. 3) and \(\psi\) the canonical mapping of \(\mathbf{Y}_\alpha \times \mathbf{Y}_\beta\) onto \(\prod_{i \in \mathbf{I}} \mathbf{Y}_i\) .

PROPOSITION 11. Let \((\mathbf{X}_t)_{t \in I}, (\mathbf{Y}_t)_{t \in I}, (\mathbf{Z}_t)_{t \in I}\) be three families of sets and let \((g_t)_{t \in I}, (g_t')_{t \in I}\) be two families of functions, all having the same index set, such that \(g_t\) is a mapping of \(\mathbf{X}_t\) into \(\mathbf{Y}_t\) and \(g_t'\) a mapping of \(\mathbf{Y}_t\) into \(\mathbf{Z}_t\) , for each \(t \in I\) . Let \(g\) and \(g'\) be the extensions of the families \((g_t)_{t \in I}\) and \((g_t')_{t \in I}\) to products. Then the extension of the family \((g_t' \circ g_t)_{t \in I}\) to products is equal to \(g'\circ g\) .

This follows immediately from Definition 2.

COROLLARY. Let \((\mathbf{X}_i)_{i\in I}, (\mathbf{Y}_i)_{i\in I}\) be two families of sets and let \((g_i)_{i\in I}\) be a family of functions. If \(g_i\) is an injection (resp. surjection) of \(\mathbf{X}_i\) into \(\mathbf{Y}_i\) for each \(i\in I\) , then the extension \(g\) of \((g_i)_{i\in I}\) to products is an injection (resp. surjection) of \(\prod_{i\in I} \mathbf{X}_i\) into \(\prod_{i\in I} \mathbf{Y}_i\) .

\begin{enumerate}
\def\labelenumi{(\arabic{enumi})}
\item
  Let us assume that \(X_{i} \neq \emptyset\) for each \(i \in I\) ; otherwise the result is trivial. Suppose that \(g_{i}\) is injective for each \(i \in I\) , and let \(r_{i}\) be a retraction of \(g_{i}\) (§3, no. 8, Definition 11), so that \(r_{i} \circ g_{i}\) is the identity mapping of \(X_{i}\) . Let r be the extension to products of the family \((r_{i})_{i \in I}\) ; since \(r \circ g\) is the extension to products of the family of identity mappings \(I_{X_{i}}\) , \(r \circ g\) is the identity mapping of \(\prod_{i \in I} X_{i}\) and hence g is injective (§3, Proposition 8).
\item
  Suppose that \(g_{i}\) is a surjection of \(X_{i}\) onto \(Y_{i}\) for each \(i \in I\) , and let \(s_{i}\) be a section of \(g_{i}\) ( \(\S 3\) , no. 8, definition 11), so that \(g_{i} \circ s_{i}\) is the identity mapping of \(Y_{i}\) . If \(s\) is the extension to products of the family \((s_{i})_{i \in I}\) , then \(g \circ s\) is the extension to products of the family of identity mappings \(I_{Y_{i}}\) and is therefore the identity mapping of \(\prod_{i \in I} Y_{i}\) ; hence \(g\) is surjective ( \(\S 3\) , Proposition 8).
\end{enumerate}

Let \((\mathbf{X}_i)_{i\in I}\) be a family of sets, and let E be a set. For every mapping \(f\) of E into \(\prod_{i\in I} \mathbf{X}_i\) , \(\operatorname{pr}_i \circ f\) is a mapping of E into \(\mathbf{X}_i\) . If \(\overline{f}\) is the extension to products of this family of mappings, and if \(d\) is the diagonal mapping of E into \(\mathbf{E}^{\mathrm{I}}\) , then it is immediate that \(f = \overline{f} \circ d\) . Conversely, let \((f_{t})_{t \in \mathbf{I}}\) be a family of functions such that \(f_{t}\) is a mapping of E into \(\mathbf{X}_{t}\) for each \(i \in \mathbf{I}\) , and let \(\overline{f}\) be the extension to products of this family; then we have \(\mathbf{pr}_{t} \circ (\overline{f} \circ d) = f_{t}\) for each \(t \in \mathbf{I}\) . By abuse of language, the mapping \(\overline{f} \circ d\) is also written as \((f_{t})_{t \in \mathbf{I}}\) . In this way we define a one-to-one mapping of the set \(\prod_{i \in \mathbf{I}} \mathbf{X}_{t}^{\mathbf{E}}\) onto the set \(\left(\prod_{i \in \mathbf{I}} \mathbf{X}_{t}\right)^{\mathbf{E}}\) ; this mapping and its inverse are said to be canonical.

